problem:  
Let \((M, \mathcal{H}, g_{\mathcal{H}})\) be a complete, equiregular sub-Riemannian manifold of Hausdorff dimension \(Q\). Denote by \(\mu\) the corresponding Hausdorff measure and by \(P_{\mathcal{H}}(E)\) the horizontal perimeter of a measurable subset \(E \subset M\). The **isoperimetric profile** of \(M\) is  
\[
I(V) = \inf\{ P_{\mathcal{H}}(E) : \mu(E) = V\}.
\]

Let \(G\) be a Lie group acting transitively on \(M\) by sub-Riemannian isometries (so that \(M \simeq G/H\) for some closed subgroup \(H\)). For a fixed basepoint \(x_0 \in M\), let \(\mathcal{G}_{x_0}\) be the nilpotentization (in the sense of Mitchell–Bellaïche) of \((M, \mathcal{H}, g_{\mathcal{H}})\) at \(x_0\), i.e., the local tangent Carnot group model.  

**Open problem (asymptotic sub-Riemannian isoperimetric rigidity):**  
Assume that \(M\) has nonnegative horizontal Ricci curvature in the sense of the generalized curvature-dimension inequality \(CD(K, Q)\) with \(K \ge 0\). Does the following limit exist and coincide with the sharp isoperimetric constant of the nilpotentization?
\[
\lim_{V \to 0^+} \frac{I(V)}{V^{(Q-1)/Q}} = C(\mathcal{G}_{x_0}) \,?
\]
Equivalently: **Does the small-volume asymptotic of the isoperimetric profile of a homogeneous sub-Riemannian manifold depend only on the local Carnot group model at a point?**  

Why is it a "good" problem:  
This problem formulates a precise conjecture linking the infinitesimal and global geometry of sub-Riemannian spaces through asymptotic isoperimetry. It asks whether the small-scale (local) isoperimetric behavior—known to be governed by the nilpotent approximation—is universal enough to determine the sharp asymptotic constant in globally homogeneous structures. The problem is grounded in foundational results (Mitchell–Bellaïche tangent cone theorem and Pansu’s isoperimetric inequality for Carnot groups) yet pushes beyond known theory into the realm of local-to-global geometric measure transitions under curvature-dimension bounds.  

Solving it would bridge three major domains:  
- **Geometric analysis**, through curvature-dimension inequalities;  
- **Sub-Riemannian geometry**, through tangent Carnot group models; and  
- **Geometric measure theory**, through perimeter and isoperimetric profiles.  

Its statement is simple and geometric, yet the analytical and structural depth is immense. A resolution could either establish a new form of asymptotic rigidity or reveal new sources of curvature-induced deviation from Carnot asymptotics, thereby reshaping the understanding of how local geometry governs global isoperimetric phenomena in sub-Riemannian spaces.